% Agent 14: Pages 396--399 (PDF pp.28--29)
% Sections 4.3 Thermal Bremsstrahlung from the Accreting Gas
%          4.4 Influence of Magnetic Fields and Synchrotron Radiation
%          4.5 Interaction of Outflowing Radiation with the Gas (beginning)
% =============================================================================

%% ----------------------------------------------------------------
%% 4.3  Thermal Bremsstrahlung from the Accreting Gas
%% ----------------------------------------------------------------

\subsection{Thermal Bremsstrahlung from the Accreting Gas}
\label{sec:4.3}

The above idealized case of spherical, adiabatic, hydrodynamic accretion can be
complicated by various physical processes.  Consider, first, the effects of energy
loss due to thermal bremsstrahlung.  To calculate the bremsstrahlung most easily,
we shall assume that the energy loss has a negligible effect on the thermal energy
density of the gas, and then we shall check that this assumption is valid.

When the effects of bremsstrahlung losses are neglected, then the temperature,
density, and velocity distributions retain their adiabatic forms (4.2.13).  One can
readily check that for $T_\infty \sim 10^4$\,K and $\Gamma = 1.4$ the gas remains nonrelativistic,
$kT \ll m_p c^2$ (but just barely so) down to $r \simeq r_s$.  Consequently, the rate at which
one gram of gas at radius $r$ radiates thermal bremsstrahlung is [cf.\ eq.~(2.1.29)]
%
\begin{equation}
\tag{4.3.1}
\varepsilon_{ff}
= \bigl(5 \times 10^{20}~\text{ergs/(g\,sec)}\bigr)\,
  (\rho / g\,\text{cm}^{-3})
  \left(\frac{T_\infty}{10^4\,\text{K}}\right)^{\!1/2}
  \left(\frac{r_i}{r}\right)^{\!-(3/a)(\Gamma-1)}
\end{equation}
%
For comparison, the rate at which adiabatic compression increases the energy of
one gram of gas is [eq.~(2.6.41)]
%
\begin{equation}
\tag{4.3.2}
\varepsilon_{\text{ad.\ heating}}
= \frac{p}{\rho}\,\frac{d\rho}{p\,dt}
= -\frac{3}{2}\,\frac{u}{r}\,\frac{p_\infty}{\rho_\infty}
  \left(\frac{r_i}{r}\right)^{-3\Gamma/2}\,.
\end{equation}
%
Using expression (4.2.12) for $r_i$, and the thermodynamic relations for a hydro\-gen gas
%
\begin{equation}
\tag{4.3.3}
a_s^2 = (\Gamma p_\infty / \rho_\infty) = \Gamma kT_\infty / m_p\,,
\end{equation}
%
we can bring this heating rate into the form
%
\begin{equation}
\tag{4.3.4}
\varepsilon_{\text{ad.\ heating}}
= \bigl(3 \times 10^4~\text{ergs}/(\text{g\,sec})\bigr)
  \left(\frac{T_\infty}{10^4\,\text{K}}\right)^{\!5/2}
  \left(\frac{M}{M_\odot}\right)^{\!-1}
  \left(\frac{r_i}{r}\right)^{-3\Gamma/2}\,.
\end{equation}

This heating rate greatly exceeds the free-free loss rate (4.3.1) at all radii.  Hence,
we are justified in our use of the adiabatic forms of $T(r)$, $\rho(r)$, and $u(r)$.  The total
power radiated as thermal bremsstrahlung is
%
\begin{equation}
\tag{4.3.5}
L_{ff}
= \int_{r_g}^{\infty} \varepsilon_{ff}\,\rho\, 4\pi r^2\, dr
\sim \left(5 \times 10^{31}~\frac{\text{ergs}}{\text{sec}}\right)
  \left(\frac{M}{M_\odot}\right)^{\!3}
  \left(\frac{T_\infty}{10^4\,\text{K}}\right)^{\!-3.3}
  \left(\frac{\rho_\infty}{10^{-24}~\text{g\,cm}^{-3}}\right)^{\!2}\,,
  \quad\text{for }\Gamma = 1.4.
\end{equation}

Perhaps a more realistic value for $\Gamma$ would be a little less than $5/3$ down to
$r \sim 10^3 r_g$ at which point $kT \sim m_e c^2$, and then $\Gamma = 4/3$ below that radius.  In
this case $L_{ff}$ would be increased by several orders of magnitude [cf.\ the rela%
tivistic corrections in eq.~(2.1.31)].  Even with such an increase, and even for
$M = 100\,M_\odot$, $L_{ff}$ is so small that it is of little or no observational interest.  For
this reason, we shall not attempt a more rigorous calculation of it.

%% ----------------------------------------------------------------
%% 4.4  Influence of Magnetic Fields and Synchrotron Radiation
%% ----------------------------------------------------------------

\subsection{Influence of Magnetic Fields and Synchrotron Radiation}
\label{sec:4.4}

Up to now we have ignored the fact that the interstellar gas possesses a magnetic
field.  For interstellar temperatures of interest, $T_\infty \sim 10^4$\,K, the magnetic field
will be frozen into the gas (cf.\ \S\ref{sec:2.5}).  The strength of the typical intergalactic
field, $B_\infty \sim 10^{-6}$\,G, is such that its energy density and pressure are not far below
those of the gas
%
\begin{equation}
\tag{4.4.1}
\frac{B_\infty^2}{8\pi}
\sim 4 \times 10^{-14}~\frac{\text{ergs}}{\text{cm}^3}\,.
\end{equation}

At large radii the field may be small enough that one can ignore its influence on
the accretion.  In this case the standard hydrodynamical inflow will stretch each
fluid element out radially (cross-sectional area of fluid element $\propto r^2$, hence
diameter $\propto r$; volume $\propto \rho^{-1} \propto r^{3/2}$, hence radial diameter $\propto r^{-1/2}$).
This radial stretch and tangential compression must quickly convert the initial
field in the fluid element into a nearly radial field with strength
%
\begin{equation}
\tag{4.4.2}
B \propto \text{(cross sectional area)}^{-1} \propto r^{-2}
\end{equation}
%
(conservation of flux).  Hence, the magnetic energy in a given fluid element will
rise as
%
\begin{equation}
\tag{4.4.3}
E_{\mathrm{mag}}
= \frac{B^2}{8\pi}~\text{(volume of element)}
\propto r^{-4}\,\rho^{-1}
\propto r^{-4}\,r^{3/2}
= r^{-5/2}\,.
\end{equation}

Not long after the gas crosses the sonic point, this magnetic energy will become
so large as to exceed the thermal energy in the fluid element
%
\begin{equation}
\tag{4.4.2a}
E_{\mathrm{therm}}
= 2^{-1}\,\frac{3\,kT}{2\,m_p}~\text{(mass of element)}
\propto r^{-(3/2)(\Gamma-1)}\,.
\end{equation}

At this point magnetic pressures must come into play, and the flow will cease
to obey the standard adiabatic hydrodynamic laws of \S\ref{sec:4.2}.

The form of the modified magnetohydrodynamic flow is not understood at
all well today.  The best guesses and calculations to date are those of Schwartzman
(1971)~\cite{Schwartzman1971}.  They suggest a turbulent flow, with reconnection of field lines between
adjacent ``cells'' (cf.\ \S\ref{sec:2.9}), and with a rough equipartition of the gravitational
potential energy between magnetic fields, kinetic infall of gas, turbulent energy,
and thermal kinetic energy of gas particles:
%
\begin{equation}
\tag{4.4.2b}
\bigl\{\text{gravitational energy}\bigr\}
= \frac{GM}{r}
\sim 4u^2
\sim 4\rho v_{\mathrm{turb}}^2
\sim 4\,\frac{3kT}{m_p}
\sim \frac{B^2}{8\pi\rho}
\end{equation}
%
per unit mass.

The mass density corresponding to this equipartition condition,
$4\pi r^2 \rho u = \dot{M}_0$, because $\dot{M}_0$ is determined by conditions
outside and at the sonic point, where the magnetic field is not yet large enough
to strongly influence the flow, we can use equation~(4.2.15b) for $\dot{M}_0$ and write
%
\begin{equation}
\tag{4.4.3a}
\rho
\sim \left(6 \times 10^{-12}~\frac{\text{g}}{\text{cm}^3}\right)
  \left(\frac{\rho_\infty}{10^{-24}~\text{g\,cm}^{-3}}\right)
  \left(\frac{T_\infty}{10^4\,\text{K}}\right)^{\!-3/2}
  \left(\frac{r_i}{r}\right)^{\!-3/2}\,.
\end{equation}

Accepting this educated guess as to the nature of the flow, we can ask what
types of radiation might be emitted.  As before~(\S\ref{sec:4.3}) bremsstrahlung is too weak
to be interesting.  However, synchrotron radiation can be rather strong.  Most of
the synchrotron radiation will come from the high-temperature, strong-field
region near the Schwarzschild radius.  There $kT \gg m_e c^2$, so the relativistic
formula~(2.3.13) for synchrotron radiation is applicable,
%
\begin{equation}
\tag{4.4.4}
\varepsilon_{\mathrm{synch}}
\simeq \bigl(2.2 \times 10^{-10}~\text{ergs}/(\text{g\,sec})\bigr)\,
  \bar{r}_e^2\, B_\perp^2\,;
\end{equation}
%
and, because the gas turns out to be optically thin, the total luminosity is
%
\begin{equation}
\tag{4.4.5}
L_{\mathrm{synch}}
\simeq \int_{r_g}^{\infty}\varepsilon_{\mathrm{synch}}\,\rho\, 4\pi r^2\, dr
\sim \bigl(1 \times 10^{32}~\text{ergs}/\text{sec}\bigr)
  \left(\frac{M}{M_\odot}\right)^{\!2}
  \left(\frac{\rho_\infty}{10^{-24}~\text{g\,cm}^{-3}}\right)^{\!2}
  \left(\frac{T_\infty}{10^4\,\text{K}}\right)^{\!-3/2}\,.
\end{equation}

For comparison, the rest mass-energy being accreted is
%
\begin{equation}
\tag{4.4.6}
\dot{M}_0 c^2
\sim \bigl(1 \times 10^{32}~\text{ergs}/\text{sec}\bigr)
  \left(\frac{M}{M_\odot}\right)^{\!2}
  \left(\frac{\rho_\infty}{10^{-24}~\text{g\,cm}^{-3}}\right)
  \left(\frac{T_\infty}{10^4\,\text{K}}\right)^{\!-3/2}\,.
\end{equation}

Since most of the radiation comes from $r \sim 2r_g$ where the thermal energy is
${\sim}\,3$ to 10 per cent of the rest mass-energy, the synchrotron losses account for
${\sim}\,3$ to 10 per cent of the thermal energy available.  However,
for $M \gtrsim 10\,M_\odot$ the effects of the energy losses on the temperature and thence
on the luminosity will not exceed a factor of ${\sim}\,1$.

Equations~(4.4.5) and (4.4.6) predict that a hole of $M \sim 10\,M_\odot$ will convert
${\sim}\,1$ per cent of the rest mass of the infalling gas into synchrotron radiation,
producing a luminosity of ${\sim}\,10^{32}$\,ergs/sec, which is about 3 per cent of the
luminosity of the sun.  The radiation, located at [cf.\ eq.~(2.3.18)] $r \simeq 2r_g$, will have a spectrum
with a broad maximum located at
%
\begin{equation}
\tag{4.4.7}
\nu_{\mathrm{peak}}
\sim (7 \times 10^{14}~\text{Hz})
  \left(\frac{\rho_\infty}{10^{-24}~\text{g\,cm}^{-3}}\right)^{\!1/2}
  \left(\frac{T_\infty}{10^4\,\text{K}}\right)^{\!-3/4}
\end{equation}
%
($\nu \sim 7 \times 10^{14}$\,Hz ($\lambda \sim 4000$\,\AA); and it will die out exponentially for
$\nu \gtrsim 7 \times 10^{14}$\,Hz; cf.\ \S\ref{sec:2.3}).  As Schwartzman (1971)~\cite{Schwartzman1971} remarks, such a spectrum
and luminosity resemble those of ``DC white-dwarf stars'' (white dwarfs with
continuous, line-free spectra).  This has led Schwartzman to speculate that
``perhaps some of the objects heretofore regarded as type~DC white dwarfs are
actually black holes.''

Instabilities and inhomogeneities in the flow (due, e.g., to reconnection of
field lines) might lead to marked fluctuations in the luminosity of the hole.  The
timescales for such fluctuations might be of the order of the gas travel time from
${\sim}\,10\,r_g$ to ${\sim}\,2r_g$; i.e.
%
\begin{equation}
\tag{4.4.8}
\Delta t_{\mathrm{fluctuations}}
\sim (10^{-3}~\text{to}~10^{-4}~\text{sec})(M/M_\odot)\,.
\end{equation}

Such rapid fluctuations in an object so faint should be very difficult to detect,
even with the 200-inch telescope.  For further details on the fluctuations, see
the end of \S\ref{sec:4.7}.

%% ----------------------------------------------------------------
%% 4.5  Interaction of Outflowing Radiation with the Gas (beginning)
%% ----------------------------------------------------------------

\subsection{Interaction of Outflowing Radiation with the Gas}
\label{sec:4.5}

As the luminosity $L$ pours out from the region $r \sim 2r_g$, it interacts with the
inflowing gas.  For the luminosities $L \sim 10^{32}$\,ergs/sec and frequencies $\nu \sim 10^{14}$\,Hz
derived in the last section, one can readily verify that the infalling gas is optically
thin to free-free absorption and to synchrotron reabsorption.  (See \S\ref{sec:2.6} for
basic concepts and equations used in such a calculation.)

What of electron scattering?  Electron scattering can exert 3 types of influence:
(i)~If the electron concentration is sufficiently high, it can act as a ``blanket'' to
impede the outflow of the radiation.  To test for such blanketing one can examine
the optical depth of the gas to electron scattering.
%
\begin{equation}
\tag{4.5.1}
\tau_{es}(r = 2r_g)
= \int_{2r_g}^{\infty}\kappa_{es}\,\rho\, dr
\simeq (1 \times 10^{-9})(M/M_\odot)
  (\rho_\infty / 10^{-24}~\text{g\,cm}^{-3})
  (T_\infty / 10^4\,\text{K})^{-3/2}\,.
\end{equation}

(A more careful calculation would replace the nonrelativistic absorption
coefficient $\kappa_{es} = 0.4$\,cm$^2$/g by the coefficient appropriate to relativistic
electrons, since $T \sim 10^{12}$\,K at $r \sim r_g$; such a calculation would also reveal
extreme optical thinness.)  (ii)~Since the outgoing photons have energies
$h\nu \sim 10$\,eV far less than the electron kinetic energies, Compton scattering can
modify the spectrum, producing high-energy photons [see eq.~(2.4.9)].  However,
the extreme optical thinness guarantees that only about one photon in a million
scatters; so this effect can be ignored.  (iii)~The outpouring photons exert a
pressure on the infalling gas when they scatter, and this effect can retard the
inflow.  Notice that the magnitude of the first two effects is governed by the
